\chapter{Bost--Connes Operators and Hochschild Trace Classes}
\label{v4-arith:connes-soibelman}

Multiplication by a prime raises its occupation number by one.
This arithmetic shift is a normalized creation operator on a bosonic
Fock space. Together with the phases indexed by
$\mathbb Q/\mathbb Z$, it gives a faithful representation of the
Bost--Connes $C^*$-algebra.
The object $\mathcal V^{\mathrm{prim}}$ is a Hochschild trace class.
Relating it to this representation requires a specified comparison map,
pairing and completion.

\section{The Bost--Connes \(C^{*}\)-Algebra}
\label{v4-arith:cs:bost-connes-explicit}

\begin{definition}[Bost--Connes presentation]
\label{v4-arith:def:bost-connes-algebra}
The Bost--Connes Q-lattice \(C^{*}\)-algebra \(C_{\mathbb{Q}}^{*}\) has
generators
\begin{align}
\{\mu_{n}\}_{n\in\mathbb{N}^{\times}},\quad
\{\mu_{n}^{*}\}_{n\in\mathbb{N}^{\times}},\quad
\{e(\gamma)\}_{\gamma\in\mathbb{Q}/\mathbb{Z}},
\end{align}
with relations
\begin{align}
\mu_1\;&=\;1,\qquad e(0)=1,\qquad e(\gamma)^*=e(-\gamma), \\
\mu_m^*\mu_n\;&=\;\mu_n\mu_m^*\quad\text{if }\gcd(m,n)=1, \\
\mu_{n}\mu_{m}\;&=\;\mu_{nm}, \\
\mu_{n}^{*}\mu_{n}\;&=\;1, \\
\mu_{n}\mu_{n}^{*}\;&=\;e_{n}\;=\;\frac{1}{n}\sum_{\gamma:n\gamma=0}e(\gamma), \\
e(\gamma_{1})e(\gamma_{2})\;&=\;e(\gamma_{1}+\gamma_{2}), \\
\mu_{n}e(\gamma)\mu_{n}^{*}\;&=\;\rho_{n}(e(\gamma))\;=\;\frac{1}{n}\sum_{\delta:n\delta=\gamma}e(\delta).
\end{align}
The \(e(\gamma)\) generate the group \(C^{*}\)-algebra
\(C^{*}(\mathbb{Q}/\mathbb{Z})\); the \(\mu_{n}\) furnish an endomorphism
semigroup action; \(\mu_{n}^{*}\mu_{n}=1\), \(\mu_{n}\mu_{n}^{*}=e_{n}\)
makes \(\mu_{n}\) a Toeplitz isometry. The crossed product
\(C^{*}(\mathbb{Q}/\mathbb{Z})\rtimes\mathbb{N}^{\times}\) yields
\(C_{\mathbb{Q}}^{*}\).
\end{definition}

\begin{theorem}[KMS states and the partition function, Bost--Connes~Prop.~3.3 and Thm~5.4]
\label{v4-arith:thm:BC-KMS}
The KMS\({}_{\beta}\) state \(\varphi_{\beta}\)
on \(C_{\mathbb{Q}}^{*}\) is unique for \(0<\beta\leq1\). For
\(\beta>1\) the KMS\({}_{\beta}\) states form a simplex with extremal
points indexed by the symmetry-broken arithmetic boundary. In the
standard Gibbs representation, for \(\beta>1\), the partition function is
\begin{equation}
Z_{\mathrm{BC}}(\beta)\;=\;\mathrm{tr}(e^{-\beta H})\;=\;\sum_{n\geq 1}n^{-\beta}\;=\;\zeta(\beta),
\end{equation}
where \(H=\log N\) generates \(\sigma_{t}(\mu_{n})=n^{it}\mu_{n}\). The
\(\zeta\)-partition function is explicitly motivated by Julia's
primon gas (\emph{Statistical theory of numbers}, Springer
Proc.~Phys.~47, 1990).
\end{theorem}

\section{Prime occupations and oscillator operators}
\label{v4-arith:cs:type-mismatch}

\begin{definition}[Prime occupation space]
\label{v4-arith:def:zero-mode-lattice}
For each prime $p$, let $\mathscr F_p=\ell^2(\mathbb N_0)$ over
$\mathbb C$, with orthonormal states $|r\rangle_p$ and vacuum $|0\rangle_p$.
Let $L_{\mathrm{alg}}$ be the complex span of tensors
$|\mathbf r\rangle=\bigotimes_p|r_p\rangle_p$ with finitely many
nonzero occupations. Declare this basis orthonormal and write $L$
for its Hilbert completion.
On $L_{\mathrm{alg}}$, define
\[
 a_p|\mathbf r\rangle
   =\sqrt{r_p}\,|\mathbf r-\mathbf e_p\rangle,\qquad
 a_p^\dagger|\mathbf r\rangle
   =\sqrt{r_p+1}\,|\mathbf r+\mathbf e_p\rangle,\qquad
 N_p=a_p^\dagger a_p,
\]
where $\mathbf e_p$ has occupation one at $p$ and zero elsewhere.
The first expression is zero when $r_p=0$.
On this common invariant dense domain,
\[
 [a_p,a_q^\dagger]=\delta_{pq}1,\qquad
 \langle a_pv,w\rangle=\langle v,a_p^\dagger w\rangle.
\]
In particular,
$|r\rangle_p=(a_p^\dagger)^r|0\rangle_p/\sqrt{r!}$.
In a vacuum Heisenberg module these are occupations of the
mode $a_{-1}^{(p)}$, rather than charges of the zero mode $a_0^{(p)}$.
\end{definition}

\begin{proposition}[Arithmetic shifts on the occupation space]
\label{v4-arith:obs:CCR-vs-Cuntz-Toeplitz}
\label{v4-arith:cor:fock-projections-determine-rep}

Let $v_p(m)$ denote the exponent of $p$ in $m$.
The map
\[
 U:\ell^2(\mathbb N^\times)\longrightarrow L,\qquad
 U\delta_m=\bigotimes_p|v_p(m)\rangle_p
\]
is unitary. Define $\mu_n\delta_m=\delta_{nm}$ on
$\ell^2(\mathbb N^\times)$.
For each prime,
\[
 U\mu_pU^{-1}=a_p^\dagger(N_p+1)^{-1/2},
\]
where the right side is first evaluated on $L_{\mathrm{alg}}$
and then extended boundedly.
For every $n\geq1$, $\mu_n^*\mu_n=1$.
The operator $\mu_n\mu_n^*$ is the orthogonal projection onto
the closed span of $\delta_m$ with $n\mid m$.
Its spectrum is $\{0,1\}$ for $n>1$ and $\{1\}$ for $n=1$.
\end{proposition}
\begin{proof}
Unique factorization gives a bijection of the two orthonormal bases,
so $U$ extends to a unitary.
The diagonal factor $(N_p+1)^{-1/2}$ cancels the creation coefficient
$\sqrt{r_p+1}$. The resulting operator raises $r_p$ by one
with coefficient one. It preserves the norm on $L_{\mathrm{alg}}$
and therefore has a unique isometric extension.
Directly,
\[
 \mu_n^*\delta_m=
 \begin{cases}\delta_{m/n},&n\mid m,\\0,&n\nmid m.\end{cases}
\]
This proves both projection identities.
For $n>1$, $\delta_n$ belongs to the range and $\delta_1$
belongs to its orthogonal complement. Both spectral values occur.
\end{proof}

\section{Arithmetic phases and the comparison carrier}
\label{v4-arith:cs:zero-mode-action}

\begin{proposition}[Bost--Connes operators on $L$]
\label{v4-arith:thm:BC-acts-on-L}

Fix the character $\chi(\gamma)=e^{2\pi i\gamma}$ of
$\mathbb Q/\mathbb Z$.
Together with the transported shifts, the operators
\[
 e(\gamma)U\delta_m=\chi(m\gamma)U\delta_m
\]
give a faithful representation of $C_{\mathbb Q}^*$ on $L$.
\end{proposition}
\begin{proof}
The phase operators are unitary and multiply by addition of their
indices. The finite character sum is
\[
 \frac1n\sum_{n\delta=\gamma}\chi(m\delta)=
 \begin{cases}
 \chi((m/n)\gamma),&n\mid m,\\
 0,&n\nmid m.
 \end{cases}
\]
These are exactly the diagonal entries of
$\mu_ne(\gamma)\mu_n^*$.
The semigroup and coprime adjoint relations hold on the occupation
basis. Thus the universal presentation defines a representation $\pi$.

For faithfulness, give
$K=\operatorname{Hom}(\mathbb Q_{>0}^{\times},U(1))$
its compact product topology, with one circle for each prime.
Its gauge action $\alpha$ fixes $e(\gamma)$ and sends $\mu_n$ to
$z(n)\mu_n$. These formulas preserve every defining relation.
The linear span of $\mu_r e(\gamma)\mu_s^*$ is dense in
$C_{\mathbb Q}^*$.
Indeed, the transfer relation implies
$e(\gamma)\mu_n=\mu_ne(n\gamma)$ and its adjoint identity.
Writing $r=dr'$, $s=ds'$ with $\gcd(r',s')=1$ gives
$\mu_r^*\mu_s=\mu_{s'}\mu_{r'}^*$.
These identities reduce products of normal terms to the same form.
A normal term has gauge degree zero exactly when $r=s$.
The transfer relation then puts it in the phase algebra.
Thus Haar averaging
$E(a)=\int_K\alpha_z(a)\,dz$ is a conditional expectation onto
the phase algebra $D$.

The universal group algebra $C^*(\mathbb Q/\mathbb Z)$ maps onto $D$.
Its representation on $\ell^2(\mathbb N^\times)$ evaluates
$C(\widehat{\mathbb Z})$ at the positive integers.
Every residue class modulo every positive integer contains a positive
integer, so these evaluation points are dense.
This representation is faithful, and hence
$D=C^*(\mathbb Q/\mathbb Z)=C(\widehat{\mathbb Z})$.

The expectation $E$ is faithful.
For $a\geq0$ and any state $\omega$, the function
$z\mapsto\omega(\alpha_z(a))$ is continuous and nonnegative.
If $E(a)=0$, its integral is zero and its value at the identity is zero.
States separate positive elements, so $a=0$.
Finally, $V_z\delta_m=z(m)\delta_m$ implements the gauge action
in $\pi$ before transport through $U$.
If $\pi(a)=0$, gauge equivariance gives $\pi(E(a^*a))=0$.
Faithfulness on $D$ gives $E(a^*a)=0$, and faithful averaging gives $a=0$.
\end{proof}

\begin{remark}[The map to a Hochschild trace class]
\label{v4-arith:obs:zero-mode-action-maximal}
The preceding construction takes place on $L$.
A Hilbert realization from Hochschild homology requires a specified
complex $(C_\bullet,b)$ over $\mathbb C$, a degree $d$, and a map
$J:L_{\mathrm{alg}}\to Z_d(C_\bullet,b)$.
It also requires a positive sesquilinear form $q$ on
$H_d(C_\bullet,b)$, a null space $\mathcal N_q$, and the completion
\[
 \mathcal H_q=\overline{H_d(C_\bullet,b)/\mathcal N_q}.
\]
The identity
$q([Jv],[Jw])=\langle v,w\rangle_L$ extends $J$ isometrically
to $L\to\mathcal H_q$.
Surjectivity onto a proposed closed sector requires density of the image.
Identifying that sector or a class in it with
$\mathcal V^{\mathrm{prim}}$ is further data.
An action on this homology and its intertwining with $J$ must also
be specified. The occupation-space construction supplies none of
these Hochschild comparison data.
\end{remark}

\begin{example}[A coherent state in the arithmetic representation]
For a prime $p$ and $z\in\mathbb C$, the vector
\[
 e^{-|z|^2/2}\sum_{k\geq0}\frac{z^k}{\sqrt{k!}}U\delta_{p^k}
\]
has norm one and lies in $L$.
It is the image of $U\delta_1$ under
$e^{-|z|^2/2}\sum_{k\geq0}z^k(U\mu_pU^{-1})^k/\sqrt{k!}$.
This operator series converges in norm, since the ratio of successive
coefficient magnitudes is $|z|/\sqrt{k+1}$.
Thus the coherent state belongs to the arithmetic algebra orbit
of the vacuum.
\end{example}

\section{The tensor decomposition by current modes}
\label{v4-arith:cs:tensor-decomposition}

\begin{proposition}[Mode-one and higher-mode factors]
\label{v4-arith:thm:BC-tensor-placement}

For a finite set $F$ of primes, let
$\mathscr F_F^{\mathrm{cur}}=\Gamma_s(\ell^2(F\times\mathbb N_{\geq1}))$
be the complex bosonic Fock space, where $\Gamma_s$ is the Hilbert
direct sum of the symmetric tensor powers.
This model has one oscillator per prime and mode.
The mode index $j\geq1$ corresponds to $a_{-j}^{(p)}/\sqrt j$.
Splitting the one-particle space into $j=1$ and $j\geq2$ gives
the canonical unitary
\[
 \mathscr F_F^{\mathrm{cur}}
 \cong L_F\widehat\otimes\mathscr F_F^{\geq2},\qquad
 L_F=\widehat\bigotimes_{p\in F}\mathscr F_p,\qquad
 \mathscr F_F^{\geq2}=\Gamma_s(\ell^2(F\times\mathbb N_{\geq2})).
\]
On algebraic number states it also preserves the grading that assigns
degree $j$ to a mode-$j$ creation operator.
Let $\mathbb N_F$ consist of the positive integers supported on $F$.
The formulas $\mu_n\delta_m=\delta_{nm}$ for $n,m\in\mathbb N_F$
and $e(\gamma)\delta_m=\chi(m\gamma)\delta_m$ act on
$L_F\cong\ell^2(\mathbb N_F)$.
They extend to the displayed tensor product by the identity on
$\mathscr F_F^{\geq2}$.
\end{proposition}
\begin{proof}
Each normalized current number state determines a unique pair of
occupation arrays, for modes one and modes at least two.
This orthonormal basis bijection gives the unitary and preserves
the sum of the mode degrees.
The subspace $\ell^2(\mathbb N_F)\subset\ell^2(\mathbb N^\times)$
is invariant under the phases and under $\mu_n,\mu_n^*$ for
$n\in\mathbb N_F$.
Restricting these operators and tensoring with the identity preserves
their relations.
\end{proof}

\begin{corollary}[Stable descendant sectors]
\label{v4-arith:cor:arithmetic-content-on-lattice}

For every closed subspace $K\subset\mathscr F_F^{\geq2}$,
the subspace $L_F\widehat\otimes K$ is invariant under the extended
arithmetic representation and its adjoints.
If $F\ne\varnothing$, the higher-mode vacuum sector is therefore
not a maximal stable subspace.
\end{corollary}
\begin{proof}
Each operator has the form $A\otimes1$ and preserves
$L_F\widehat\otimes K$ together with its adjoint.
When $F\ne\varnothing$, the higher-mode space contains both the
vacuum and a nonzero mode-two state. The vacuum sector is properly
contained in the full stable tensor product.
\end{proof}

\section{Kesten--McKay = Plancherel of Unramified \(GL_{2}(\mathbb{Q}_{p})\)}
\label{v4-arith:cs:kesten-mckay-plancherel}

\begin{theorem}[Kesten 1959; McKay 1981; Lubotzky 1994~Thm~7.4.5]
\label{v4-arith:thm:kesten-mckay-plancherel}
The Kesten--McKay distribution on
\([-2\sqrt{p},2\sqrt{p}]\) for \(q=p\) prime,
\begin{equation}
d\nu_{p}(x)\;=\;\frac{(p+1)\sqrt{4p-x^{2}}}{2\pi((p+1)^{2}-x^{2})}\,dx,
\end{equation}
is the spectral measure of the adjacency operator on the
\((p+1)\)-regular tree, equivalently the \(p\)-adic Bruhat--Tits tree
of \(\mathrm{PGL}_{2}(\mathbb{Q}_{p})\), equivalently the Plancherel
measure of the spherical-function algebra on the tree.
\end{theorem}

\begin{corollary}
\label{v4-arith:cor:kesten-mckay-is-local-hecke}
The Kesten--McKay averaging is the
Plancherel decomposition of the \(GL_{2}(\mathbb{Q}_{p})\) spherical
Hecke algebra, not a spectral theorem for the \(GL_{1}\) Bost--Connes
semigroup algebra. Any map from the Julia--Bost--Connes partition
function to \(\nu_{p}\) must pass through a specified \(GL_{1}\to GL_{2}\)
correspondence.
\end{corollary}

\begin{remark}[arithmetic averaging and Plancherel comparison]
\label{v4-arith:rem:arithmetic-averaging-plancherel}
The finite occupation-phase operator \(A_S\) of
Definition~\ref{v4-arith:def:averaging-naive} acts on a polynomial oscillator
space. It does not construct the proposed arithmetic carrier or
quotient for \(\mathrm{av}^{\mathrm{arith}}\).
An arithmetic Plancherel interpretation requires a specified
representation, its spectral transform, and an intertwining map from
the intended arithmetic object. The operadic averaging
\(\mathfrak g^{E_1}\to\mathfrak g^{\mathrm{mod}}\) of Vol~I instead
acts on chiral cohomology. A common integral form does not identify
these carriers or supply a map between them.
\end{remark}

\section{The Soibelman Cohomological Hall Algebra Reading}
\label{v4-arith:cs:soibelman-coha}

A rank-\(n\) representation of the Jordan quiver over
\(\mathbb F_q\) is an endomorphism of \(\mathbb F_q^n\).
Change of basis acts by conjugation. Its finite groupoid mass weights
each isomorphism class by the reciprocal of its automorphism-group order.

\begin{proposition}[finite-field mass of Jordan-quiver representations]
\label{v4-arith:thm:KS-CoHA-adelic}
For a prime power \(q\) and \(n\geq1\),
\[
 \sum_{[A]\subset M_n(\mathbb F_q)}
 \frac{1}{|\operatorname{Cent}_{GL_n(\mathbb F_q)}(A)|}
 =\frac{q^{n^2}}{|GL_n(\mathbb F_q)|}
 =\prod_{i=1}^n(1-q^{-i})^{-1}.
\]
The sum is over conjugacy classes of matrices. In rank one its value
is \(q/(q-1)\).
\end{proposition}
\begin{proof}
An orbit of \(A\) has
\(|GL_n(\mathbb F_q)|/|\operatorname{Cent}_{GL_n(\mathbb F_q)}(A)|\)
elements. Summing orbit sizes gives \(|M_n(\mathbb F_q)|=q^{n^2}\).
Counting successive independent columns gives
\[
 |GL_n(\mathbb F_q)|=\prod_{j=0}^{n-1}(q^n-q^j)
 =q^{n^2}\prod_{i=1}^n(1-q^{-i}).
\]
Division proves the formula and its rank-one specialization.
\end{proof}

This scalar groupoid mass is not a real spectral measure.
A comparison with the Kesten--McKay measure requires a specified
operator, state and moment or spectral-transform map. A Behrend-weighted
or vanishing-cycle refinement also requires its own normalization.
Kontsevich--Soibelman, arXiv:1006.2706, \S\S2.4--2.5, constructs the
cohomological Hall product and its shuffle expression. That construction
does not supply the spectral comparison described here.

\begin{proposition}[three-way comparison obstruction]
\label{v4-arith:prop:HY-BC-Soibelman-dictionary}
The composed arrow
\begin{equation*}
\text{Julia primon gas}
\xrightarrow{\text{Plancherel}}
\mathrm{Bost\text{-}Connes\,Hecke\text{-}Plancherel}
\xrightarrow{\text{CoHA-motivic}}
\mathrm{Soibelman\,DT\,on}\,\mathcal{Q}_{\mathrm{Jordan}}
\end{equation*}
is not defined by the preceding theorems. The Bost--Connes algebra is
a \(GL_{1}\) semigroup system, the Kesten--McKay measure is a
\(GL_{2}\) spherical Plancherel measure, and the Jordan-quiver CoHA is
a motivic Hall algebra. Equality of Euler factors can be asserted only
after a functor between these three structures has been constructed.
\end{proposition}

\section{Gauge Derivations and Smooth Domains}
\label{v4-arith:cs:HH-comparison}

The occupation representation supplies bounded arithmetic shifts and
phases. Its time derivative is unbounded, so a Hochschild calculation
must specify an algebra on which that derivative is defined. We use
the finite normal-term algebra to construct gauge classes and their
cup products. Smoothness under finite prime-circle actions gives a
continuous version; smoothness under the time evolution gives a
different coefficient algebra. Their distinction is visible on an
explicit norm-convergent series. None of these constructions identifies
an arithmetic Hochschild trace carrier.

Put \(A=C_{\mathbb Q}^*\) over \(\mathbb C\), represented faithfully
on \(\mathcal L=\ell^2(\mathbb N^\times)\) as in
Proposition~\ref{v4-arith:thm:BC-acts-on-L}. Write its orthonormal
basis as \(e_m\), so that
\(\mu_n e_m=e_{nm}\) and
\(e(\gamma)e_m=\exp(2\pi i m\gamma)e_m\).
Define the complex vector space of finite sums
\[
 A_{\mathrm{alg}}=
 \operatorname{span}_{\mathbb C}
 \{\mu_r e(\gamma)\mu_s^*:r,s\ge1,\ \gamma\in\mathbb Q/\mathbb Z\}.
\]
It is a dense unital \(*\)-subalgebra of \(A\). Indeed, if
\(d=\gcd(s,u)\), \(s=ds'\), and \(u=du'\), the defining relations give
\begin{align*}
 &(\mu_r e(\gamma)\mu_s^*)(\mu_u e(\eta)\mu_v^*)\\
 &\hspace{2em}=\mu_{ru'}e(u'\gamma+s'\eta)\mu_{vs'}^*.
\end{align*}
Here \(\mu_s^*\mu_u=\mu_{u'}\mu_{s'}^*\),
\(e(\gamma)\mu_{u'}=\mu_{u'}e(u'\gamma)\), and
\(\mu_{s'}^*e(\eta)=e(s'\eta)\mu_{s'}^*\).
Adjoint exchanges \(r,s\) and negates \(\gamma\).
Thus the displayed span is closed under multiplication and adjoint,
and it contains all generators.

\begin{theorem}[gauge classes in algebraic Hochschild cohomology]
\label{v4-arith:thm:CM-HH-BC}
Use algebraic \(\mathbb C\)-multilinear Hochschild cochains of
\(A_{\mathrm{alg}}\), with coefficients in its regular bimodule.
For every \(\lambda=(\lambda_p)_{p\ \mathrm{prime}}
\in\prod_p\mathbb C\), the formula
\[
 D_\lambda(\mu_r e(\gamma)\mu_s^*)
 =\left(\sum_p(v_p(r)-v_p(s))\lambda_p\right)
       \mu_r e(\gamma)\mu_s^*
\]
defines a derivation. There is an injective linear map
\[
 \prod_p\mathbb C\longhookrightarrow
 HH^1_{\mathrm{alg}}(A_{\mathrm{alg}},A_{\mathrm{alg}}),
 \qquad \lambda\longmapsto[D_\lambda].
\]
No boundedness or summability condition on the family is required.
Every nonzero \(D_\lambda\) is unbounded for the inherited
\(C^*\)-norm and cannot be implemented by a bounded operator on
\(\mathcal L\). Moreover
\(HH^0_{\mathrm{alg}}(A_{\mathrm{alg}},A_{\mathrm{alg}})=\mathbb C\).
This is a lower bound in degree one, not a complete Hochschild calculation.
\end{theorem}
\begin{proof}
Let \(K=\prod_p U(1)\), and for \(z\in K\) put
\(z(n)=\prod_p z_p^{v_p(n)}\).
The assignments \(\alpha_z(\mu_n)=z(n)\mu_n\) and
\(\alpha_z(e(\gamma))=e(\gamma)\) preserve the relations and define
an action by \(*\)-automorphisms. It is norm continuous in \(z\)
on each finite normal sum, and hence on \(A\) by uniform
approximation and isometry. A normal term has character \(z(r/s)\).
Haar integration against the inverse character projects onto its
homogeneous component; distinct rational degrees therefore give
\[
 A_{\mathrm{alg}}=\bigoplus_{g\in\mathbb Q_{>0}^{\times}}
 (A_{\mathrm{alg}})_g,
 \qquad (A_{\mathrm{alg}})_g(A_{\mathrm{alg}})_h
 \subseteq(A_{\mathrm{alg}})_{gh}.
\]
This does not assert independence of normal terms of the same degree.
Unique factorization makes
\(\lambda(g)=\sum_pv_p(g)\lambda_p\) an additive character of the
multiplicative rational group: \(\lambda(gh)=\lambda(g)+\lambda(h)\).
All these sums are finite. Define \(D_\lambda\) on a homogeneous
element of degree \(g\) by multiplication by \(\lambda(g)\).
The direct-sum decomposition proves well-definedness, and the character
identity proves the Leibniz rule. In particular it annihilates the
phases and sends \(\mu_n\) to \(\lambda(n)\mu_n\) and
\(\mu_n^*\) to \(-\lambda(n)\mu_n^*\).

Algebraic Hochschild one-cocycles are precisely derivations and
one-coboundaries are inner derivations. If \(D_\lambda=[b,-]\)
for a bounded operator \(b\), the isometry \(\mu_p^k\) gives
\[
 k|\lambda_p|=\|D_\lambda(\mu_p^k)\|
 =\|[b,\mu_p^k]\|\le2\|b\|\qquad(k\ge1).
\]
Hence every \(\lambda_p=0\). Applying this to differences proves
injectivity on Hochschild classes. The same estimate on the
norm-one sequence \(\mu_p^k\) proves unboundedness whenever a
prime value is nonzero, and excludes a bounded extension \(A\to A\).

An operator commuting with all phases has vanishing matrix entries
between distinct \(e_m,e_n\): some \(\gamma\in\mathbb Q/\mathbb Z\)
has different characters at \(m,n\). It is therefore diagonal,
say \(Te_m=t_me_m\). Commutation with \(\mu_n\) gives
\(t_{nm}=t_m\); setting \(m=1\) makes all entries equal.
Thus the represented commutant is scalar. Faithfulness implies
\(Z(A_{\mathrm{alg}})=\mathbb C\), which is its Hochschild degree zero.
\end{proof}

The product \(\prod_p\mathbb C\) parametrizes actual derivations,
not convergent infinite operator sums: every input has finite
homogeneous support. A derivation in this family preserves adjoint
exactly when every \(\lambda_p\) is purely imaginary, as follows
by comparing its values on \(\mu_p\) and \(\mu_p^*\).
The preceding algebraic result does not compute continuous
Hochschild cohomology of the Banach algebra \(A\).

\subsection{Higher Gauge Cup Products}

The degree-one classes also detect higher cohomology. To prove this,
we use a scalar character to change coefficients and then evaluate
cup products on explicit finite cycles. This coefficient change is
part of the argument; it is not an identification with Hochschild
homology with regular coefficients.

There is a character \(\varepsilon:A\to\mathbb C\) defined by
\[
 \varepsilon(\mu_n)=\varepsilon(\mu_n^*)=
 \varepsilon(e(\gamma))=1.
\]
These values respect the unit, involution, phase multiplication,
semigroup, coprime-adjoint and isometry relations. The transfer
relation becomes one equal to the average of \(n\) ones.
The universal property therefore gives a unital \(*\)-homomorphism,
of norm one. It is not the vacuum vector state: the latter takes
\(\mu_p\) to zero.

Let \(\mathbb C_\varepsilon\) be the bimodule with
\(a\cdot c\cdot b=\varepsilon(a)c\varepsilon(b)\).
Our algebraic cochains have differential
\begin{align*}
 (bF)(a_1,\ldots,a_{n+1})={}&a_1F(a_2,\ldots,a_{n+1})\\
 &+\sum_{i=1}^n(-1)^iF(a_1,\ldots,a_ia_{i+1},\ldots,a_{n+1})\\
 &+(-1)^{n+1}F(a_1,\ldots,a_n)a_{n+1}.
\end{align*}
Composition with \(\varepsilon\) is a cochain map from regular
coefficients to \(\mathbb C_\varepsilon\), by multiplicativity.
For scalar-coefficient chains, written \(c[a_1|\cdots|a_n]\), use
\begin{align*}
 b(c[a_1|\cdots|a_n])={}&c\varepsilon(a_1)[a_2|\cdots|a_n]\\
 &+\sum_{i=1}^{n-1}(-1)^ic[a_1|\cdots|a_ia_{i+1}|\cdots|a_n]\\
 &+(-1)^n\varepsilon(a_n)c[a_1|\cdots|a_{n-1}].
\end{align*}
Evaluation makes these two differential formulas dual.

\begin{lemma}[finite detecting cycles]
\label{v4-arith:bc:gauge-cycles}
For distinct primes \(P=(p_1,\ldots,p_n)\), put \(s_j=\mu_{p_j}\).
Then
\[
 c_P=\sum_{\sigma\in S_n}\operatorname{sgn}(\sigma)
           1[s_{\sigma(1)}|\cdots|s_{\sigma(n)}]
\]
is a cycle with coefficients in \(\mathbb C_\varepsilon\).
Consequently the functional on regular-coefficient cochains
\[
 L_P(F)=\sum_{\sigma\in S_n}\operatorname{sgn}(\sigma)
 \varepsilon\bigl(F(s_{\sigma(1)},\ldots,s_{\sigma(n)})\bigr)
\]
annihilates Hochschild coboundaries.
\end{lemma}
\begin{proof}
For each interior face, pair a permutation with the one exchanging
the adjacent multiplied factors. Their signs are opposite and the
boundary terms agree, since the \(s_j\) commute. Pair the first face
of \((j_1,\ldots,j_n)\) with the last face of its cyclic rotation
\((j_2,\ldots,j_n,j_1)\). Rotation has sign \((-1)^{n-1}\),
whereas the last boundary sign is \((-1)^n\); the terms cancel
because \(\varepsilon(s_j)=1\). For \(n=1\), the boundary is
\(\varepsilon(s_1)-\varepsilon(s_1)=0\).
The cochain map and the dual differential formulas now show that
evaluation on \(c_P\) vanishes on every regular-coefficient coboundary.
\end{proof}

\begin{theorem}[an algebraic exterior-algebra injection]
\label{v4-arith:bc:exterior-injection}
Let \(W=\operatorname{Hom}_{\mathbb Z}
 (\mathbb Q_{>0}^{\times},(\mathbb C,+))=\prod_p\mathbb C\).
The gauge derivations give an injective graded-algebra map
\[
 \Lambda_{\mathbb C}^*W\longhookrightarrow
 HH^*_{\mathrm{alg}}(A_{\mathrm{alg}},A_{\mathrm{alg}}),
 \qquad \lambda\longmapsto[D_\lambda].
\]
The exterior algebra here is algebraic: its elements are finite
sums of finite wedges, without completion. In particular the
Hochschild cohomology is nonzero in every nonnegative degree.
\end{theorem}
\begin{proof}
The cup of \(n\) derivations is
\[
 F_{\lambda_1,\ldots,\lambda_n}(a_1,\ldots,a_n)
   =D_{\lambda_1}(a_1)\cdots D_{\lambda_n}(a_n).
\]
Substitution into the differential gives
\(b(F\smile G)=bF\smile G+(-1)^{\deg F}F\smile bG\): the
interior terms cancel except the two indicated coboundaries.
Each derivation is a cocycle, so their cup is a cocycle.
Using \(D_{\lambda_i}(s_j)=\lambda_i(p_j)s_j\) gives
\[
 L_P(F_{\lambda_1,\ldots,\lambda_n})
       =\det\bigl(\lambda_i(p_j)\bigr)_{1\le i,j\le n}.
\]
For two derivations \(D,E\), their composition is a one-cochain.
Expanding both Leibniz rules yields
\[
 b(DE)(a,b)=-D(a)E(b)-E(a)D(b).
\]
Thus their cohomology classes anticommute, and each square vanishes
in characteristic zero. This defines the exterior-algebra map.

An exterior element involves a finite-dimensional subspace
\(U\subset W\). Prime-coordinate functionals separate its points;
their span is \(U^*\), since otherwise their common annihilator
would contain a nonzero element with all prime coordinates zero.
A finite subset therefore spans \(U^*\), and wedges of its
coordinates span \((\Lambda^nU)^*\). Every nonzero degree-\(n\)
exterior element pairs nontrivially with some wedge of distinct
prime coordinates. The corresponding determinant functional \(L_P\)
detects its image in Hochschild cohomology. Degree zero is detected
by \(\varepsilon(1)=1\).
\end{proof}

\begin{example}[degrees two and three]
For distinct primes \(p,q\), the chain
\(1[\mu_p|\mu_q]-1[\mu_q|\mu_p]\) has boundary
\[
 1[\mu_q]-1[\mu_p\mu_q]+1[\mu_p]
 -1[\mu_p]+1[\mu_q\mu_p]-1[\mu_q]=0.
\]
Here \(D_p\) denotes the family with value one at \(p\) and zero
elsewhere. The evaluation
\(L_{(p,q)}(D_p\smile D_q)=1\) proves this cup is not a coboundary.
For three distinct primes, among the six terms defining
\(L_{(p,q,r)}(D_p\smile D_q\smile D_r)\), only the identity
assignment is nonzero, and its value is one. These are explicit
higher classes, not a classification of all Hochschild classes.
\end{example}

The time evolution is \(\sigma_t(\mu_n)=n^{it}\mu_n\), fixing the
phases. Its norm-generator domain consists of those \(a\in A\) for
which \(\lim_{t\to0}(\sigma_t(a)-a)/t\) exists in norm; call that
generator \(\delta\), and write \(A_\sigma^\infty\) for its
norm-smooth orbit algebra. Its closedness and smooth topology are
proved below.

\subsection{Finite-Prime and All-Prime Gauge Smoothness}

For a finite prime set \(F\), let \(\mathbb T^F\subset K\) act by
setting the other coordinates to one. Define
\[
 B_F=\{a\in A:z\mapsto\alpha_z(a)
                   \text{ is norm-smooth on }\mathbb T^F\}.
\]
For angular differentiation in the \(p\)-coordinate write
\(\delta_p\); thus \(\delta_p(\mu_q)=i\delta_{pq}\mu_q\) and
\(D_p=\delta_p/i\). For nonnegative multiindices \(\alpha\) supported
in \(F\), put \(q_\alpha(a)=\|\delta^\alpha a\|\), including
\(q_0(a)=\|a\|\). Define also
\[
 B_{\mathrm g}=\bigcap_{F\ \mathrm{finite}}B_F
\]
with all finite-support multiindex seminorms. This definition uses
ordinary smoothness on each finite-dimensional torus; it makes no
unspecified infinite-dimensional differentiability assumption.

\begin{proposition}[smooth gauge algebras and their cup classes]
\label{v4-arith:bc:gauge-smooth}
The algebras \(B_F\) and \(B_{\mathrm g}\) are Fr\'echet
\(*\)-algebras, dense in \(A\) for its \(C^*\)-norm.
The derivations \(D_p\), for \(p\in F\) on \(B_F\) and all primes
on \(B_{\mathrm g}\), are continuous. With jointly continuous
complex multilinear cochains and regular coefficients, there are
injections
\begin{align*}
 \Lambda_{\mathbb C}^*\operatorname{span}_{\mathbb C}\{D_p:p\in F\}
  &\longhookrightarrow HH^*_{\mathrm{cont}}(B_F,B_F),\\
 \Lambda_{\mathbb C}^*\left(\bigoplus_p\mathbb C D_p\right)
  &\longhookrightarrow HH^*_{\mathrm{cont}}(B_{\mathrm g},B_{\mathrm g}).
\end{align*}
\end{proposition}
\begin{proof}
If \(a_j\) is Cauchy for every \(q_\alpha\) on \(B_F\), its
orbit maps \(f_j(z)=\alpha_z(a_j)\) are Cauchy in every uniform
derivative norm, since
\(\partial^\alpha f_j(z)=\alpha_z(\delta^\alpha a_j)\)
and the action is isometric. The space of smooth Banach-valued
functions on the torus is complete: the fundamental theorem of
calculus in coordinate charts shows that uniform limits of functions
and their first derivatives retain those derivatives, and iteration
gives the assertion at every order. The limit \(f\) has
\(f(z)=\alpha_z(a)\), where \(a=\lim_j a_j\) in norm, by continuity
of each automorphism. Thus \(a\in B_F\) and all seminorms converge.
The iterated product rule
\[
 \delta^\alpha(ab)=\sum_{\beta\le\alpha}\binom\alpha\beta
       (\delta^\beta a)(\delta^{\alpha-\beta}b)
\]
proves continuous multiplication. Adjoint commutes with real angular
derivatives and is continuous; mixed derivatives commute by smoothness.

There are countably many finite prime sets and finite-support
multiindices. A sequence Cauchy in all seminorms of \(B_{\mathrm g}\)
converges in every \(B_F\); the limits coincide with its norm limit
in \(A\). This proves completeness of \(B_{\mathrm g}\), and the
same product rule proves its Fr\'echet algebra structure.
Finite normal sums belong to every \(B_F\), giving norm density.
For \(p\) and finite \(F\), differentiate the orbit for
\(\mathbb T^{F\cup\{p\}}\) and restrict to \(\mathbb T^F\).
This shows that \(\delta_p\) preserves \(B_{\mathrm g}\), with
\(q_\alpha(\delta_p a)=q_{\alpha+e_p}(a)\); the assertion on
\(B_F\) is identical for \(p\in F\).

The finite cup cochains are continuous by these estimates and
continuity of multiplication. The character is continuous because
\(|\varepsilon(a)|\le q_0(a)\). The finite cycles of
Lemma~\ref{v4-arith:bc:gauge-cycles} therefore annihilate every
continuous coboundary as well. Their determinant evaluations prove
the two injections. No completed chain-complex comparison is needed.
\end{proof}

\begin{proposition}[which algebraic gauge derivations extend continuously]
\label{v4-arith:bc:gauge-continuity}
An algebraic gauge derivation \(D_\lambda\) extends to a continuous
linear map \(B_{\mathrm g}\to A\) exactly when \(\lambda\) has
finite support. In that case it extends as a continuous derivation
\(B_{\mathrm g}\to B_{\mathrm g}\). For \(B_F\), the analogous
criterion is that \(\lambda\) be supported in \(F\).
\end{proposition}
\begin{proof}
Continuity of an extension \(T:B_{\mathrm g}\to A\) gives
\[
 \|T(a)\|\le C\max_{\alpha\in E}q_\alpha(a)
\]
for a finite set \(E\) of multiindices, enlarged to include zero.
Let \(F\) contain their supports. If a prime \(q\notin F\), then
\(q_0(\mu_q^k)=1\), whereas every positive derivative in \(E\)
is zero. Thus \(k|\lambda_q|\le C\) for every \(k\), forcing
\(\lambda_q=0\). Conversely a finite linear combination of the
continuous \(D_p\) is continuous into \(B_{\mathrm g}\).
The same argument on \(B_F\) proves the second criterion.
These criteria concern this family of algebraic derivations;
they do not classify all continuous derivations.
\end{proof}

\begin{example}[gauge smoothness does not imply time differentiability]
Enumerate distinct primes \(p_j\), and put
\[
 a=\sum_{j\ge1}2^{-j}\mu_{p_j}^{2^j}.
\]
The series converges in norm. For each finite \(F\), only finitely
many summands move under \(\mathbb T^F\), while the remaining
norm-convergent sum is fixed. Hence \(a\in B_{\mathrm g}\).
If its time derivative existed in norm, its value on \(e_1\)
would have coordinate \(i\log p_j\) at each distinct state
\(e_{p_j^{2^j}}\). These coordinates are not square summable.
Thus \(a\notin\operatorname{Dom}(\delta)\), and
\(B_{\mathrm g}\) is not contained in \(A_\sigma^\infty\).
The time family \(\lambda_p=i\log p\) also fails the finite-support
continuity criterion. Gauge-smooth and time-smooth Hochschild
complexes are therefore different specified constructions.
\end{example}

\begin{proposition}[the closed time generator and its proper domain]
\label{v4-arith:obs:CM-stricter-than-H-arith}
For \(\sigma_t(\mu_n)=n^{it}\mu_n\) and
\(\sigma_t(e(\gamma))=e(\gamma)\), let
\[
 \operatorname{Dom}(\delta)=
 \left\{a\in A:\lim_{t\to0}\frac{\sigma_t(a)-a}{t}
                    \text{ exists in }A\right\}.
\]
The norm generator \(\delta\) is closed and densely defined, its
domain is proper, and its restriction to \(A_{\mathrm{alg}}\) is
\(D_{(i\log p)_p}\).
\end{proposition}
\begin{proof}
The gauge action makes \(\sigma\) pointwise norm continuous.
Every normal term is entire analytic, with derivative
\(i\log(r/s)\mu_r e(\gamma)\mu_s^*\); hence the domain contains
the dense normal algebra. Fix a prime \(p\). The series
\[
 a=\sum_{j\ge1}2^{-j}\mu_p^{2^j}
\]
converges in norm because the summand norms sum to one. Applied
to \(e_1\), its orbit difference quotient has coordinate
\[
 2^{-j}\frac{\exp(it2^j\log p)-1}{t}
 \longrightarrow i\log p
\]
at \(e_{p^{2^j}}\). If the quotient converged in operator norm,
its value on \(e_1\) would converge to an \(\ell^2\) vector having
this same nonzero coordinate at infinitely many orthonormal states.
This is impossible, so \(a\notin\operatorname{Dom}(\delta)\).

For a domain element, differentiation along the norm-continuous
orbit gives the Bochner integral identity
\[
 \sigma_t(a)-a=\int_0^t\sigma_s(\delta a)\,ds.
\]
If \(a_n\to a\) and \(\delta a_n\to b\) in norm, isometry permits
passage to the limit in this identity. Dividing by \(t\) and using
norm continuity of \(\sigma_s(b)\) gives \(\delta a=b\).
This proves closedness.
\end{proof}

\begin{proposition}[a smooth Hochschild class]
\label{v4-arith:bc:smooth-HH1}
The algebra of norm-smooth elements is
\[
 A_\sigma^\infty=\bigcap_{j\ge0}\operatorname{Dom}(\delta^j),
 \qquad p_k(a)=\sum_{j=0}^k\|\delta^j a\|.
\]
With these seminorms it is a Fr\'echet algebra, dense in \(A\).
The derivation \(\delta:A_\sigma^\infty\to A_\sigma^\infty\)
is continuous and represents a nonzero class in continuous
Hochschild degree one with coefficients in this same smooth algebra.
\end{proposition}
\begin{proof}
Repeated differentiation of the orbit identifies smoothness with
membership in every generator domain. If a sequence is Cauchy for
every \(p_k\), each \(\delta^j a_n\) has a norm limit. Closedness
of \(\delta\), applied successively, identifies these limits as
\(\delta^j a\). Thus the seminorm topology is complete. The
iterated Leibniz rule
\[
 \delta^j(ab)=\sum_{i=0}^j\binom ji(\delta^i a)(\delta^{j-i}b)
\]
proves closure under multiplication and its continuity; the generator
also commutes with adjoint. The normal algebra consists of smooth
elements and is norm dense in \(A\). Finally
\(p_k(\delta a)\le p_{k+1}(a)\) proves continuity. An inner
implementer in this smooth algebra would be a bounded element of
\(A\), contrary to the power estimate in
Theorem~\ref{v4-arith:thm:CM-HH-BC} with \(\lambda_p=i\log p\).
\end{proof}

\begin{proposition}[the occupation Hamiltonian and represented closure]
\label{v4-arith:bc:H-domain-and-closure}
On \(\mathcal L\) let \(He_m=(\log m)e_m\), with maximal domain
\[
 \operatorname{Dom}(H)=
 \left\{\xi\in\mathcal L:\sum_m(\log m)^2|\xi_m|^2<\infty\right\}.
\]
This diagonal operator is self-adjoint, with finite-support core.
Normal terms preserve its domain, and there
\[
 [H,\mu_r e(\gamma)\mu_s^*]
 =\log(r/s)\mu_r e(\gamma)\mu_s^*.
\]
The unitaries \(e^{itH}\) implement \(\sigma_t\), and the represented
von Neumann closure of \(A\) is \(B(\mathcal L)\).
\end{proposition}
\begin{proof}
For a real diagonal multiplier, testing the adjoint against basis
vectors gives precisely the stated square-summability domain;
thus the maximal operator equals its adjoint. Truncating a vector
in both the norm and weighted norm proves the core assertion.
A normal term vanishes unless \(s\mid m\), and otherwise sends
\(e_m\) to a phase times \(e_{rm/s}\). The logarithmic weights
change by the constant \(\log(r/s)\), which proves domain preservation
and the commutator identity. Equivalently, prove the identity on
the core and pass to graph limits by closedness. Exponentiating the
diagonal weights gives the asserted implementation. The scalar
commutant computed in Theorem~\ref{v4-arith:thm:CM-HH-BC} and the
bicommutant theorem give the represented closure \(B(\mathcal L)\).
\end{proof}

Thus \(\delta=i[H,-]\) on \(A_{\mathrm{alg}}\) has an unbounded
spatial implementer. The prime-power norm estimate in
Theorem~\ref{v4-arith:thm:CM-HH-BC} excludes every bounded implementer,
so \(\delta\) is outer on \(A_{\mathrm{alg}}\) and
\(A_\sigma^\infty\). This occupation
closure is a type-I-infinity factor. A factor obtained from another
state by a KMS--GNS construction is a different carrier. Its type
identifies none of the algebraic, Banach or Fr\'echet Hochschild
complexes above. The algebraic and all-prime gauge-smooth lower bounds above
are nonzero in every degree. They do not classify the full cohomology
on these carriers or compute the other Hochschild theories.

\begin{question}[the arithmetic Hochschild comparison]
\label{v4-arith:thm:H-arith-corrected}
Construct the arithmetic category or algebra whose Hochschild
invariants are to realize \(\mathcal V^{\mathrm{prim}}\), specifying
its coefficient ring, bimodule or categorical trace, topology,
grading and completion. Construct the comparison with the
Bost--Connes normal algebra or a specified smooth algebra and the
cycle pairing of Remark~\ref{v4-arith:obs:zero-mode-action-maximal}.
Determine which gauge classes survive under that comparison and
whether an integral degree-one carrier generated by logarithmic
prime data exists. Any proposed degree-two relation to a Witt
cotangent complex requires its actual map and coefficient comparison.
A trace class alone defines no Hochschild cochain algebra, and the
complex injection above establishes neither an integral lattice nor
concentration of the arithmetic theory in degrees zero through two.
\end{question}

\section{Vertex-Algebraic Invariants of the Heisenberg Covers}
\label{v4-arith:cs:new-content}

\begin{theorem}[universal Heisenberg-cover invariants]
\label{v4-arith:thm:universal-non-arithmetic}
The finite Heisenberg cover
\(\mathcal{H}^{\otimes \pi(N)}\) of
\(\mathcal{V}^{\mathrm{prim}}_{N}\) is a rank-\(\pi(N)\) Heisenberg
vertex algebra. It lies in the Heisenberg archetype \(\mathsf{G}\) of
the Vol~I landscape census, with cover central charge
\(\kappa=\pi(N)\), depth \(r_{\max}=2\), and chiral Hochschild
cohomology concentrated in degrees \(\{0,1,2\}\). These are invariants
of the cover, not of the trace class.
\end{theorem}

\begin{corollary}[placement of the finite covers]
\label{v4-arith:cor:placement-G-archetype}
The finite Heisenberg covers of
\(\mathcal{V}^{\mathrm{prim}}\) lie in archetype \(\mathsf{G}\)
(Heisenberg / Gaussian), with \(\kappa_{N}=\pi(N)\). The Hilbert comparison with the Hochschild trace class
requires the map, pairing and completion of
Remark~\ref{v4-arith:obs:zero-mode-action-maximal}.
\end{corollary}

\section{The Fock representation and the trace comparison}
\label{v4-arith:cs:correct-identification}

\begin{corollary}[Finite current Fock representation]
\label{v4-arith:thm:CS-corrected-placement}

The finite current Fock space has the Hilbert tensor decomposition
of Proposition~\ref{v4-arith:thm:BC-tensor-placement}.
Its arithmetic shifts and phases act on the mode-one factor and
extend by the identity on the higher-mode factor.
\end{corollary}
\begin{proof}
Apply Proposition~\ref{v4-arith:thm:BC-tensor-placement} to the given finite
set of primes.
\end{proof}

\begin{remark}[Scope of the module structure]
The tensor decomposition is a statement about current Fock spaces.
It induces no tensor decomposition of a Hochschild trace class
without a compatible comparison.
The higher-mode factor alone is not the full vacuum Heisenberg
vertex algebra: it omits all mode-one creations.
Nor does the finite-prime restriction assert faithfulness for all
arithmetic phases. If $q\notin F$, the nonzero Bost--Connes projection
$e_q=q^{-1}\sum_{j=0}^{q-1}e(j/q)$ acts by zero on $L_F$,
since $q$ divides no element of $\mathbb N_F$.
\end{remark}

\begin{remark}[the \(\mathsf{B}\)-row threshold]
A proposed arithmetic chiral algebra should encode Hecke eigenvalues
in its stress tensor. Comparing its chiral Hochschild cohomology with
Bost--Connes cohomology requires a specified algebraic or smooth
coefficient algebra and an actual subquotient map. The gauge classes
above do not construct this map.
Placement in the Borcherds/BKM \(\mathsf{B}\)-row further requires
verification of the hypotheses of the Vol~III classification theorem
(Theorem~C, BKM column), whose corresponding value is \(K^\kappa=8\).
The proposed passage from the Heisenberg-cover \(\mathsf{G}\)-row
uses a Borcherds-product encoding of Hecke data in the stress tensor,
with the target relation \(\Phi_N=c_N(0)/2\). Prime indexing alone
supplies neither that encoding nor the Hochschild comparison.
\end{remark}

\section{Structural Conclusions}
\label{v4-arith:cs:summary}

\begin{enumerate}
\item Bost--Connes \(C^{*}_{\mathbb{Q}}\) carries \(\zeta(\beta)\) as
Gibbs partition function for \(\beta>1\)
(Theorem~\ref{v4-arith:thm:BC-KMS}).
\item The oscillator operators and bounded arithmetic shifts act
on the same prime occupation space. The shift is normalized
creation, as in
Proposition~\ref{v4-arith:cor:fock-projections-determine-rep}.
\item The arithmetic phases and shifts give the faithful
representation of Proposition~\ref{v4-arith:thm:BC-acts-on-L}.
Its identification with a Hochschild trace class requires the
map, pairing and completion of
Remark~\ref{v4-arith:obs:zero-mode-action-maximal}.
\item The current Fock space splits into mode-one and higher-mode
factors. The arithmetic representation extends over the latter,
with the stable sectors of
Corollary~\ref{v4-arith:cor:arithmetic-content-on-lattice}.

\item The Kesten--McKay measure describes the spherical
\(GL_2(\mathbb Q_p)\) Hecke spectrum. Relating it to the
\(GL_1\) Bost--Connes system or the Jordan-quiver CoHA requires
the comparison maps described in
Corollary~\ref{v4-arith:cor:kesten-mckay-is-local-hecke} and
Proposition~\ref{v4-arith:prop:HY-BC-Soibelman-dictionary}.
\item The normal-term algebra has the independent gauge classes of
Theorem~\ref{v4-arith:thm:CM-HH-BC}. The time derivative requires
its specified algebraic or smooth domain; the algebraic and all-prime gauge-smooth cup products are nonzero
in every degree, without a full Hochschild classification.
\item The finite Heisenberg covers lie in archetype \(\mathsf{G}\)
with \(\kappa=\pi(N)\), \(r_{\max}=2\); the \(\mathsf{B}\)-row threshold
requires Borcherds-product encoding of Hecke data
(Corollary~\ref{v4-arith:cor:placement-G-archetype}).
\end{enumerate}

The vertex algebra structure is universal and non-arithmetic because
it belongs only to the finite Heisenberg covers. The Bost--Connes representation gives the arithmetic shifts, phases
and Gibbs partition function. Its relation to the Hochschild trace
class requires the comparison of
Remark~\ref{v4-arith:obs:zero-mode-action-maximal}. A genuine VA-arithmetic object encoding Hecke
eigenvalues in the stress tensor would be a new object, not
\(\mathcal{V}^{\mathrm{prim}}\).
